\documentclass[12pt]{article}
\usepackage[utf8]{inputenc}
\usepackage[T1]{fontenc}
\usepackage{hyperref}
\title{Walla Walla Language Resources for K-12 Curriculum Development}
\author{Compiled by Claude Sonnet 5 (claude-sonnet-5) for Andr\'as Kornai}
\date{}
\begin{document}
\maketitle

\section{Introduction}
Walla Walla (ISO 639-3 waa; Glottocode wall1238) is a Sahaptian language, treated as the northeastern dialect of Sahaptin. Searches found no standalone Walla Walla teaching materials. Most materials are for the closely related Umatilla and Yakama dialects. Sources disagree on endangerment status. This report is short because little Walla Walla-specific material turned up.

\section{Primary Educational Resources}
No Walla Walla primers, textbooks or school curricula turned up in the searches. The Confederated Tribes of the Umatilla Indian Reservation (CTUIR) language program is the likeliest source of classroom materials. It teaches Umatilla, Walla Walla and Cayuse: \url{https://ctuir.org/about/language/}

\section{Linguistic and Reference Materials}
\begin{itemize}
\item Sahaptin language overview (Wikipedia), which lists Rude's \emph{Umatilla Dictionary} (University of Washington Press, 2014) among its references: \url{https://en.wikipedia.org/wiki/Sahaptin_language}
\item Glottolog entry for Walla Walla: \url{https://glottolog.org/resource/languoid/id/wall1238}
\item Smithsonian finding aid, MS 1001, a comparative vocabulary that includes Walla Walla: \url{https://sirismm.si.edu/EADpdfs/NAA.MS1001.pdf}
\item American Philosophical Society, Sahaptin materials (Phillips Fund Collection): \url{https://indigenousguide.amphilsoc.org/entry/40100}
\end{itemize}
These are archival or related-dialect materials, not Walla Walla textbooks. I could not confirm the contents of the online Umatilla dictionary mentioned by CTUIR.

\section{Religious and Community Resources}
CTUIR is the community organisation behind the language program: \url{https://ctuir.org/about/language/}

\section{Digital and Audio Resources}
None specific to Walla Walla were found in the search results.

\section{Curriculum Development Context}
CTUIR says speakers of all levels speak Umatilla and Walla Walla, so the language is not dormant. A curriculum would most realistically be built with CTUIR, adapting Umatilla-based materials where the dialects overlap. This needs confirmation with the tribe.
\end{document}